% SGA 3，第 VII_B 篇，英文重建稿——第 1 轮。
% 范围：权威底本分篇第 51--52 页，第 2.4 节全文，含其中的定理和推论，止于 2.4.A 之前。
% 合并阅读本页码： 564--565.
% 再版印刷页码： 553--554.
% 本范围内原版页边定位： 520.
% 本范围内编者注： 90--91.
% 权威底本： Polo--Gille Exp7B-13oct24.pdf,
% SHA-256 D303690E0195654BBE7C3DAA5D89EA23E24CD709A3575217234AC19AE8836694.
% 曾参考 Jacob Reinhold 的 CC BY 4.0 Markdown 作为比对谱系。
% 原英文 Loop 1 阶段说明：三幅图采用直接取自 PDF 的精确源裁片，普通数学均为原生 LaTeX。当前中文工程已采用原生图。
\subsection*{2.4}\phantomsection \label{sga3:VIIB:2.4}\addcontentsline{toc}{subsection}{2.4. 形式群的商}

现在回到任意伪紧环 \(k\)，以便将 \SGARefLink{sga3:VIIB:1.4.thm}{1.4}--\SGARefLink{sga3:VIIB:1.5}{1.5} 中关于对拓扑平坦等价关系取商的结果应用于形式群。\footnote{\label{sga3:VIIB:nde:90}编者补充了下文原论证的细节。}

  % SGA3-VII-DIAG-106
  % SGA3-VII-DIAG-107
设 \(u:H\to G\) 为形式 \(k\)-群的单态射，\(\mu:G\times G\to G\) 为 \(G\) 的乘法态射，\(\lambda\) 为复合\begin{center}\phantomsection \label{sga3:VIIB:2.4:lambda-composite} % SGA3-VII-DIAG-106
% Authority: Exp7B-13oct24.pdf, local p. 51.
% Crop: figures/Exp7B_localp051_quotient_action_lambda_composite.png
% Locator: tex/components/42_expose_VIIB_section24_formal_group_quotients_en.tex, Loop 1 line 30.
\(
\lambda:\quad
G\times H\xrightarrow{\ \operatorname{id}_G\times u\ }
G\times G\xrightarrow{\ \mu\ }G.
\)
\end{center}。由于 \(u\) 为单态射，对\begin{center}\phantomsection \label{sga3:VIIB:2.4:equivalence-pair} % SGA3-VII-DIAG-107
% Authority: Exp7B-13oct24.pdf, local p. 51.
% Crop: figures/Exp7B_localp051_quotient_equivalence_pair.png
% Locator: tex/components/42_expose_VIIB_section24_formal_group_quotients_en.tex, Loop 1 line 37.
\begin{tikzpicture}[baseline=(lhs.base)]
  \node (lhs) {$G\times H$};
  \node[right=2.5cm of lhs] (rhs) {$G$};
  \draw[-{Straight Barb[length=4pt,width=4pt]}]
    ([yshift=2.5pt]lhs.east) --
    node[above] {$\operatorname{pr}_1$}
    ([yshift=2.5pt]rhs.west);
  \draw[-{Straight Barb[length=4pt,width=4pt]}]
    ([yshift=-2.5pt]lhs.east) --
    node[below] {$\lambda$}
    ([yshift=-2.5pt]rhs.west);
\end{tikzpicture}
\end{center}是 \(\operatorname{Vaf}/k\) 中的等价对 \textup{（参见 V，2.b）}。回顾 \textup{（参见 \SGARefLink{sga3:VIIB:1.2.C}{1.2.C}）}，此对的余核 \(G/H\) 定义如下。

设 \(\mathcal O(G)\) 和 \(\mathcal O(H)\) 分别为 \(G\) 和 \(H\) 的仿射代数，\[
  \Delta:\mathcal O(G)
  \longrightarrow
  \mathcal O(G)\widehat\otimes_k\mathcal O(G)
\] 为 \(\mathcal O(G)\) 的对角态射，\(I\) 为 \[
  u^\natural:\mathcal O(G)\longrightarrow\mathcal O(H).
\] 的核。由 \SGARefLink{sga3:VIIB:1.3.prop}{命题 1.3}，态射 \(u^\natural\) 诱导同构 \[
  \mathcal O(G)/I\xrightarrow{\ \sim\ }\mathcal O(H).
  % SGA3-VII-DIAG-108
\]。于是 \(G/H\) 的仿射代数 \(\mathcal O(G/H)\) 为态射对\begin{center}\phantomsection \label{sga3:VIIB:2.4:affine-equalizer-pair} % SGA3-VII-DIAG-108
% Authority: Exp7B-13oct24.pdf, local p. 51.
% Crop: figures/Exp7B_localp051_quotient_affine_equalizer_pair.png
% Locator: tex/components/42_expose_VIIB_section24_formal_group_quotients_en.tex, Loop 1 line 66.
\begin{tikzpicture}[baseline=(lhs.base)]
  \node (lhs) {$\mathcal O(G)$};
  \node[right=3cm of lhs] (rhs)
    {$\mathcal O(G)\widehat\otimes_k\mathcal O(H),$};
  \draw[-{Straight Barb[length=4pt,width=4pt]}]
    ([yshift=2.5pt]lhs.east) --
    node[above] {$\tau_1$}
    ([yshift=2.5pt]rhs.west);
  \draw[-{Straight Barb[length=4pt,width=4pt]}]
    ([yshift=-2.5pt]lhs.east) --
    node[below] {$(\operatorname{id}\widehat\otimes u^\natural)\Delta$}
    ([yshift=-2.5pt]rhs.west);
\end{tikzpicture}
\end{center}的核，其中 \(\tau_1(x)=x\widehat\otimes1\)。故 \[
  \mathcal O(G/H)
  =
  \bigl\{
    x\in\mathcal O(G)
    \mathrel{\big|}
    \Delta(x)-x\widehat\otimes1
    \in\mathcal O(G)\widehat\otimes I
  \bigr\}.
\]。

若进一步 \(H\) 在 \(k\) 上拓扑平坦，则 \(\operatorname{pr}_1\) 拓扑平坦，由 \SGARefLink{sga3:VIIB:1.4.thm}{定理 1.4} 得到下述结果。

\medskip\noindent\textbf{定理。}\phantomsection \label{sga3:VIIB:2.4:theorem}{\itshape 设 \(u:H\to G\) 为形式 \(k\)-群的单态射。假设 \(H\) 在 \(k\) 上拓扑平坦。则投影 \[
  p:G\longrightarrow G/H
\] 满射且拓扑平坦，存在同构 \[
  G\times H
  \xrightarrow{\ \sim\ }
  G\times_{G/H}G,
  \tag{*}\label{sga3:VIIB:2.4:equation:quotient-fiber-product-star}
\]，且 \(G/H\) 表示平坦拓扑的商层。

因此 \(G/H\) 具有由 \(G\) 作用的对象的典范结构，且 \(p:G\to G/H\) 为等变态射。若进一步 \(u\) 将 \(H\) 等同于 \(G\) 的正规子群，则 \(G/H\) 具有典范的形式 \(k\)-群结构，使 \(p:G\to G/H\) 为形式 \(k\)-群态射，且 \(H\) 为 \(p\) 的核。}

事实上，第一个断言由 \SGARefLink{sga3:VIIB:1.4.thm}{1.4} 得出，另两个由 \SGARefLink{sga3:iv:corollary:5:2:2}{第四讲， 推论 5.2.2} 和 \SGARefLink{sga3:VIIB:5.2.4}{5.2.4} 得出。

\medskip\noindent\textbf{推论。}\phantomsection \label{sga3:VIIB:2.4:corollary}\footnote{\label{sga3:VIIB:nde:91}编者将此推论明确写出；例如 \SGARefLink{sga3:VIIB:5.2.1:section-proposition}{5.2.1} 和 \SGARefLink{sga3:VIIB:5.2.3}{5.2.3} 中要用到它。}{\itshape 设 \(G\) 为形式 \(k\)-群，\(H\) 为 \(G\) 的形式子群。设 \(A\)、\(A/J\) 和 \(B\) 分别为 \(G\)、\(H\) 和 \(G/H\) 的仿射代数；\(I_A\) 为 \(A\) 的增广理想，并置 \(I_B=B\cap I_A\)。假设 \(H\) 在 \(k\) 上拓扑平坦。则 \[
  J=AI_B,
\]，即由 \(I_B\) 生成的闭理想。}

事实上，投影 \(B\to B/I_B\) 对应于 \(G/H\) 的单位截面 \[
  e:\operatorname{Spf}(k)\longrightarrow G/H
\]。由 \textup{(*)}，\(H\) 与纤维积 \[
  \operatorname{Spf}(k)\times_{G/H}G.
\] 等同。因此其仿射代数 \(A/J\) 与 \[
  (B/I_B)\widehat\otimes_BA
  \simeq
  A/AI_B.
\] 等同。
